\subsection{Sum-Product Algorithm}
The Sum-Product Algorithm (also called belief propagation) operates by iteratively exchanging messages between variable and check nodes.
Messages are expressed as LLRs, representing beliefs about the value of each bit.

Initialization:
each variable node $v_i$ sends its channel LLR $\lambda_i$ to all connected check nodes.

Iterative Message Passing:
Let:
\begin{itemize}
    \item $m_{i \to j}$: message from variable node $i$ to check node $j$.
    \item $m_{j \to i}$: message from check node $j$ to variable node $i$.
    \item $\mathbb{N}_j$: neighbors of node $i$.
\end{itemize}

(a) Variable Node Update:
Each variable node sends to check node $j$:
$$m_{i \to j} = \lambda_i + \sum_{j' \in \mathcal{N}(i) \setminus j} m_{j' \to i}$$

Here, the node aggregates all messages from neighboring check nodes except $j$, combining them with the channel information.

(b) Check Node Update:
Each check node sends to variable node $i$:
$$m_{j \to i} = 2 \tanh^{-1} \left( \prod_{i' \in \mathcal{N}(j) \setminus i} \tanh\left( \frac{m_{i' \to j}}{2} \right) \right)$$

This step enforces parity constraints by combining the incoming beliefs from all other connected variable nodes. 
The hyperbolic tangent function emerges from the exact inference in tree-structured graphs and operates in the LLR domain.

This uses the identity:
$$\tanh^{-1}(x) = \frac{1}{2} \log \left( \frac{1 + x}{1 - x} \right)$$

Final Decision:
After a fixed number of iterations or convergence, the final LLR for each variable node is computed as:
$$L_i = \lambda_i + \sum_{j \in \mathcal{N}(i)} m_{j \to i}$$

A hard decision is made based on the sign of $L_i$:
$$
\hat{x}_i =
\begin{cases}
0 & \text{if } L_i \geq 0 \\
1 & \text{if } L_i < 0
\end{cases}
$$

Finally, the estimated codeword $\hat{x}$ is checked against the parity-check condition:
$$H \cdot \hat{x}^T = 0 (mod 2)$$

If this condition is satisfied, decoding is successful \cite{910572}.